\chapter{直线与多边形}\label{ch:g4:geometry}

用直尺、三角尺和圆规，几何就成了一门手艺。
本章从工具操作的角度，介绍直线的两种特殊位置——
\omterm{def:g4:geometry:def}{垂直}与\omterm{def:g4:geometry:def}{平行}，并带你认识\omterm{def:g4:geometry:polygon}{多边形}家族。
（这些概念的理论见
\cref{ch:g6:lines}）。

\section{垂直与平行}

\begin{definition}[垂直、平行]\label{def:g4:geometry:def}
两条直线相交成\omterm{def:g3:shapes:rightangle}{直角}时，称它们\emph{垂直}\index{垂直}（用三角尺检验，
并用小方框标记）。
若两条直线\emph{平行}\index{平行}，即它们方向相同且永不相交——
就像铁路上的两根钢轨，
始终保持相同的距离。
\end{definition}

\begin{method}[用三角尺作图]\label{met:g4:geometry:draw}
\emph{作直线 $d$ 的垂线，过点 $P$}：将\omterm{def:g3:shapes:rightangle}{直角}的一边
沿 $d$ 放置，滑动三角尺，使另一边经过 $P$，
然后画线。

\emph{作 $d$ 的\omterm{def:g4:geometry:def}{平行}线，过点 $P$}：先将三角尺的一边
沿 $d$ 放置，再将直尺紧靠三角尺的另一边作导轨，
沿直尺滑动三角尺至 $P$，画线。
\end{method}

\begin{omfigure}
\begin{tikzpicture}[scale=0.85]
  \draw[thick] (-0.4,0) -- (4.2,0) node[below, font=\small] {$d$};
  \draw[omDef, thick] (1.8,-0.9) -- (1.8,2.3);
  \draw[thin] (1.8,0) ++(-0.25,0) -- ++(0,0.25) -- ++(0.25,0);
  \fill (1.8,1.6) circle (1.6pt) node[right, font=\small] {$P$};
  \node[below, font=\small] at (1.9,-1.1) {过 $P$ 的垂线};
  \begin{scope}[xshift=7.4cm]
  \draw[thick] (-0.4,0) -- (4.2,0) node[below, font=\small] {$d$};
  \draw[omThm, thick] (-0.4,1.6) -- (4.2,1.6);
  \fill (1.6,1.6) circle (1.6pt) node[above, font=\small] {$P$};
  \draw[Stealth-Stealth, thin] (3.4,0) -- (3.4,1.6)
    node[midway, right, font=\small] {各处间距相同};
  \node[below, font=\small] at (1.9,-1.1) {过 $P$ 的平行线};
  \end{scope}
\end{tikzpicture}

{\small 两种作图方法。\omterm{def:g4:geometry:def}{平行}线之间的间距保持不变：
在两处量取间距是一种很好的检验方法。}
\end{omfigure}

\begin{example}[在图形中辨认垂直与平行]\label{ex:g4:geometry:spot}
在长方形中，每条边都与其相邻的两条边\omterm{def:g4:geometry:def}{垂直}，
并与对边\omterm{def:g4:geometry:def}{平行}。方格纸就是一张现成的\omterm{def:g4:geometry:def}{平行}线与垂线网格，
这正是它特别适合作图的原因，
用起来十分顺手。
\end{example}

\section{多边形}

\begin{definition}[多边形]\label{def:g4:geometry:polygon}
边界由线段组成的封闭图形称为\emph{多边形}\index{多边形}，
各线段称为它的\emph{边}；角上的点称为
\emph{顶点}\index{顶点}。多边形按边数命名：
\emph{三角形}（$3$）、\emph{四边形}\index{四边形}（$4$），
\emph{五边形}\index{五边形}（$5$）、\emph{六边形}\index{六边形}（$6$），
\emph{八边形}\index{八边形}
（$8$）。
\end{definition}

\begin{omfigure}
\begin{tikzpicture}[scale=0.62]
  \draw[thick, fill=omDef!10] (0,0) -- (2,0) -- (1.2,1.7) -- cycle;
  \node[below, font=\small] at (1,-0.3) {三角形};
  \begin{scope}[xshift=3.4cm]
  \draw[thick, fill=omThm!10] (0,0) -- (2.2,0) -- (2.5,1.5) --
    (0.5,1.8) -- cycle;
  \node[below, font=\small] at (1.25,-0.3) {四边形};
  \end{scope}
  \begin{scope}[xshift=7.4cm]
  \draw[thick, fill=omMeth!10] (90:1.1) -- (162:1.1) -- (234:1.1) --
    (306:1.1) -- (18:1.1) -- cycle;
  \node[below, font=\small] at (0,-1.4) {五边形};
  \end{scope}
  \begin{scope}[xshift=10.9cm]
  \draw[thick, fill=omProp!10] (0:1.1) -- (60:1.1) -- (120:1.1) --
    (180:1.1) -- (240:1.1) -- (300:1.1) -- cycle;
  \node[below, font=\small] at (0,-1.4) {六边形};
  \end{scope}
\end{tikzpicture}

{\small \omterm{def:g4:geometry:polygon}{多边形}家族（图中画出的\omterm{def:g4:geometry:polygon}{五边形}和\omterm{def:g4:geometry:polygon}{六边形}都是
\emph{正\omterm{def:g4:geometry:polygon}{多边形}}：各边相等，各角也相等）。圆不是\omterm{def:g4:geometry:polygon}{多边形}，
其边界不由线段组成。}
\end{omfigure}

\begin{example}[描述多边形]\label{ex:g4:geometry:describe}
一个有趣的游戏：把图形描述得足够准确，
让同学不看图也能画出来。“一个\omterm{def:g4:geometry:polygon}{四边形} $ABCD$，其中
$AB = 5$ cm，$A$ 处是\omterm{def:g3:shapes:rightangle}{直角}，$AD = 3$ cm\,\dots”——
准确的词汇是让信息传达成功的关键。
\end{example}

\section{圆与圆规}

\begin{definition}[圆]\label{def:g4:geometry:circle}
圆心为 $O$、\omterm{def:g3:shapes:shapes}{半径}为 $3$ cm 的\emph{圆}是所有
到定点的距离恰为 $3$ cm 的点的集合（定点即 $O$）。圆规画圆时，针尖放在
$O$ 处，张开 $3$ cm，旋转一周即可。\emph{直径}\index{直径}是
经过圆心、连接圆上两点的线段，其长度为
\omterm{def:g3:shapes:shapes}{半径}的两倍。
\end{definition}

\begin{example}[圆规：保持距离的工具]\label{ex:g4:geometry:compass}
圆规不止能画圆，它还能\emph{搬运
距离}。要标出所有与某点的距离恰为 $4$ cm 的点（该点记为 $A$），画出
圆心为 $A$、\omterm{def:g3:shapes:shapes}{半径}为 $4$ cm 的圆；要比较两条线段而不用直尺，
只需在其中一条上张开圆规，再到另一条上检验。
\end{example}

\begin{method}[在方格纸上复制图形]\label{met:g4:geometry:reproduce}
\begin{enumerate}
  \item 在空白方格纸上选定一个起始\omterm{def:g4:geometry:polygon}{顶点}及其位置；
  \item 沿图形逐个\omterm{def:g4:geometry:polygon}{顶点}移动，数出方格的
        步数（“右 $3$ 格、上 $2$ 格”）；
  \item 依次连接各\omterm{def:g4:geometry:polygon}{顶点}，并与原图比较：
        边长相等，方向一致。
\end{enumerate}
\end{method}

\section{练习}

\begin{exercise}[$\star$]\label{exo:g4:geometry:1}
画一条直线 $d$ 和不在其上的一点 $P$。先作
与 $d$ \omterm{def:g4:geometry:def}{垂直}且经过 $P$ 的直线，再作与 $d$ \omterm{def:g4:geometry:def}{平行}
且经过 $P$ 的直线（用直尺平移法）。标出\omterm{def:g3:shapes:rightangle}{直角}。
\end{exercise}

\begin{exercise}[$\star$]\label{exo:g4:geometry:2}
观察大写字母 E、H、N、T、Z。每个字母中，找出一对
\omterm{def:g4:geometry:def}{平行}的笔画和一对\omterm{def:g4:geometry:def}{垂直}的笔画，
若存在则指出。
\end{exercise}

\begin{exercise}[$\star$]\label{exo:g4:geometry:3}
说出下列\omterm{def:g4:geometry:polygon}{多边形}的名称：$6$ 条边；$4$ 条边；$3$ 条边；$8$ 条边；
$5$ 条边。每个\omterm{def:g4:geometry:polygon}{多边形}有几个\emph{\omterm{def:g4:geometry:polygon}{顶点}}？
\end{exercise}

\begin{exercise}[$\star$]\label{exo:g4:geometry:4}
画一个任意\omterm{def:g4:geometry:polygon}{五边形}（不必是正\omterm{def:g4:geometry:polygon}{五边形}），并标出它的\omterm{def:g4:geometry:polygon}{顶点}
$A$、$B$、$C$、$D$、$E$。它有几条边？画出从 $A$ 出发的
对角线（连接 $A$ 与不相邻\omterm{def:g4:geometry:polygon}{顶点}的线段）：
共有几条？
\end{exercise}

\begin{exercise}[$\star$]\label{exo:g4:geometry:5}
画一个圆心为 $O$、\omterm{def:g3:shapes:shapes}{半径}为 $4$ cm 的圆；画一条\omterm{def:g4:geometry:circle}{直径} $[AB]$；
并在圆上取一点 $M$，它不在 $[AB]$ 上。直接说出 $OA$，
$OM$ 和 $AB$ 的长度（不要测量）。
\end{exercise}

\begin{exercise}[$\star$]\label{exo:g4:geometry:6}
利用圆规这一保持距离的工具
（见\cref{ex:g4:geometry:compass}），比较同学所画三角形的各边，
并判断它是否有两条边相等。
\end{exercise}

\begin{exercise}[$\star$]\label{exo:g4:geometry:7}
在方格纸上画出下面的路线，并将其围成一个\omterm{def:g4:geometry:polygon}{多边形}：
从点 $A$ 出发；向右 $4$ 格，向上 $2$ 格，向左 $3$ 格，向上 $2$ 格；
再沿直线回到 $A$。所得\omterm{def:g4:geometry:polygon}{多边形}叫什么名称
（数一数它有几条边）？
\end{exercise}

\begin{exercise}[$\star$]\label{exo:g4:geometry:8}
判断下列说法是否正确，并画图说明理由。
“两条\omterm{def:g4:geometry:def}{垂直}的直线总是相交。”
“两条\omterm{def:g4:geometry:def}{平行}的直线永不相交。”
“一条直线可以既与一条直线\omterm{def:g4:geometry:def}{平行}，又与
另一条直线\omterm{def:g4:geometry:def}{垂直}。”
\end{exercise}

\begin{exercise}[$\star\star$]\label{exo:g4:geometry:9}
画一条直线 $d$，再画一条直线 $e$，使其与 $d$ \omterm{def:g4:geometry:def}{垂直}，然后画
直线 $f$，使其与 $e$ \omterm{def:g4:geometry:def}{垂直}。观察 $d$ 与 $f$，你发现了什么？
这一观察将在\cref{ch:g6:lines}中成为一个定理。
\end{exercise}

\begin{exercise}[$\star\star$]\label{exo:g4:geometry:10}
不让同学看图，用文字向他描述下面的图形：
一个边长为 $4$ cm 的正方形，其中有一个\omterm{def:g3:shapes:shapes}{半径}为 $2$ cm 的圆，
圆心在正方形的中心（与四条边都相切）。写出
这段描述，并写成编号的步骤。
\end{exercise}

\begin{exercise}[$\star\star$]\label{exo:g4:geometry:11}
一个\omterm{def:g4:geometry:polygon}{四边形}有几条对角线？\omterm{def:g4:geometry:polygon}{五边形}呢？\omterm{def:g4:geometry:polygon}{六边形}呢？
画出这三种情形，仔细数一数，并描述你
所发现的规律。
\end{exercise}
